\documentclass[11pt]{report} % Clase de documento para un informe
% Documento y codificación general
\usepackage[utf8]{inputenc} % Tildes, "ñ" y símbolos especiales.
% Idioma y traducción
\usepackage[spanish,es-tabla]{babel} % Traduce el documento al español.
\usepackage[a4paper,left=3cm,right=2.5cm,top=2.5cm,bottom=2.5cm]{geometry} % Permite cambiar el tamaño de la página.
% Imágenes y figuras
\usepackage{graphicx} % Permite incluir imágenes.
\usepackage{subcaption} % Permite agregar subfiguras dentro de una figura principal, cada una con su leyenda.
\usepackage{float} % Permite fijar la posición de las figuras y tablas [H].
% Matemáticas y símbolos
\usepackage{amsmath} % Permite usar entornos matemáticos avanzados.
\usepackage{amssymb} % Permite usar símbolos matemáticos adicionales.
\usepackage{mathtools} % Permite usar herramientas matemáticas avanzadas.
% Hipervinculos y referencias
\usepackage[hidelinks]{hyperref} % Permite crear hipervínculos transparentes dentro del documento.
\usepackage[table,xcdraw]{xcolor} % Permite usar colores textos y tablas.
\usepackage{setspace} % Permite ajustar el interlineado.


\begin{document}
    \begin{spacing}{1.5}
        \begin{flushleft}

                \begin{center}
                    \Large
                    \textbf{FORMULARIO DE AVANCE MENSUAL DE\\
                    PROYECTO INTEGRADOR}
                \end{center}
                \vspace{2em}

                FECHA: 07/04/2026

                ALUMNO: Tardón, Damián David

                TEMA: Diseño, desarrollo e implementación de un software para la adquisición y análisis de ondas de impulso tipo rayo (1,2/50 \(\mu\)s).

                DIRECTOR: Blasco, Marcos Javier

                CO-DIRECTOR: Serra, Gabriel Horacio.

                FECHA ESTIMADA DE FINALIZACIÓN: Mayo 2026.

                \hrulefill

                INFORME:

        \end{flushleft}
    \end{spacing}

    Se optimizó el sistema de gestión de archivos, permitiendo la navegación a través del sistema de directorios.\par
    Se actualizó la interfaz gráfica (GUI) mediante la incorporación de parámetros de ensayo adicionales y nuevas funcionalidades orientadas a mejorar la usabilidad.\par
    Se integraron los módulos de procesamiento digital de señales, el control del osciloscopio y la interfaz de usuario en único archivo ejecutable (.exe), permitiendo la adquisición y el análisis de ondas de impulso de manera centralizada.\par
    Se realizaron ensayos con señales reales en colaboración con el personal del Laboratorio de Alta Tensión, obteniendo resultados satisfactorios en la detección y caracterización de impulsos individuales, aunque se detectaron dificultades técnicas al procesar secuencias de ondas consecutivas.\par
    Actualmente, se trabaja en el ajuste de los umbrales de tolerancia para garantizar una discriminación correcta entre variaciones normales y fallas reales, ya que la sensibilidad excesiva ante pequeñas desviaciones en la forma de onda genera actualmente falsos positivos (detección errónea de descargas disruptivas).\par
    Las actividades descritas en este informe se corresponden con las tareas 6 y 7 del diagrama de Gantt propuesto en la SAT.\par

\end{document}